% Source: https://artofproblemsolving.com/community/c3h2064226_functional_equations
% Stable identity: aops-topic-2064226-q01
\emph{The zero function.} If $f(x_0)=0$ for some $x_0$, then $f(x)=f(x-x_0)f(x_0)=0$ for all $x$. So $f\equiv0$ is one solution.

\emph{Otherwise $f>0$.} If $f$ never vanishes, then $f(x)=f\left(\dfrac {x}{2}\right)^2>0$ for all $x$. Let $g(x)=\ln f(x)$; then $g$ is continuous and
\[
g(x+y)=g(x)+g(y).
\]
Let $c=g(1)$. Then $g(0)=2g(0)$ gives $g(0)=0$, and $g(-x)=-g(x)$. By induction $g(nx)=ng(x)$ for positive integers $n$, hence for all integers. For a rational $\dfrac {p}{q}$ ($q>0$), $q\,g\left(\dfrac {p}{q}\right)=g(p)=pc$, so $g\left(\dfrac {p}{q}\right)=\dfrac {p}{q}\,c$. Every real $x$ is a limit of rationals, and $g$ is continuous, so $g(x)=cx$ for all $x$. Hence $f(x)=e^{cx}=a^x$ with $a=e^c>0$.

Conversely, $f\equiv0$ and $f(x)=a^x$ ($a>0$) satisfy the equation. The solutions are
\[
\boxed{f\equiv0\quad\text{or}\quad f(x)=a^x\ \text{ for some constant }a>0}.
\]
(The case $a=1$ gives the constant $1$. Without continuity there are also wild non-measurable solutions.)
